\documentclass{librairies/lib}
\usepackage{glossaries}

%%%%%%%%%%%%%VARIABLE EN TETE%%%%%%%%%%%
\newcommand{\TITRE}{\today}
\newcommand{\NOM}{CHAUSSON Thibault}
\newcommand{\UVXX}{AV00}
%%%%%%%%%%%%%%%%%%%%%%%%%%%%%%%%%%



\makeglossaries


\newglossaryentry{ENSEEIHT}
{
    name=ENSEEIHT,
    description={École Nationale Supérieure d'Électrotechnique, d'Électronique, d'Informatique, d'Hydraulique et des Télécommunications}
}

\newglossaryentry{CNJE}
{
    name=CNJE,
    description={Confédération nationale des Junior-Entreprises}
}


\begin{document}

%!TEX root = ../main.tex

\thispagestyle{empty}
    \begin{sffamily}
        \begin{center}

            % Upper part of the page. The '~' is needed because \\
            % only works if a paragraph has started.

            %\includegraphics[scale=0.2]{ressources/logo}~\\[1cm]


            % Upper part of the page. The '~' is needed because \\
            % only works if a paragraph has started.

            % \includegraphics[scale=0.2]{ressources/logo}~\\[1cm]

            \begin{figure}[htbp]
                \centering
                \begin{minipage}[c]{0.45\textwidth}
                    \centering
                    \includegraphics[width=\textwidth]{ressources/logo}
                    \label{fig:image1}
                \end{minipage}
                \hfill
                \begin{minipage}[c]{0.45\textwidth}
                    \centering
                    \includegraphics[width=\textwidth]{ressources/AE couleur}
                    \label{fig:image2}
                \end{minipage}
            \end{figure}

            \vspace{1cm}

            \textsc{\LARGE Université de technologie de Belfort-Montbéliard}\\[2cm]

            \textsc{\Large Réalisation audiovisuelle}\\[1cm]

            % Title
            {\color{UTBMcolor}\HRule} \\[0.4cm]
            { \huge \bfseries Promotion de la vie associative\\[0.4cm] }
            {\color{UTBMcolor}\HRule} \\[1,5cm]


            % Author and supervisor

            \begin{center}
                Réalisé par : CHAUSSON Thibault
    
                \vspace{0,5cm}
                
                Encadré par : MIRABEL-GRAFF Corinne \& JOUFFROY François
            \end{center}




            \vfill

            % Bottom of the page
            \textsc{\large \today}

        \end{center}
    \end{sffamily}




%\maketitle
\newpage
\pagestyle{no_number}

\section*{Préambule}
\addcontentsline{toc}{section}{Préambule}

J'ai la volonté d'écrire l'ensemble de se rapport en utilisant l'écriture inclusive et plus particulièrement le point médian\footnote{Unicode : U+00B7}.
L'ensemble des règles appliquées sont issues du Manuel d’Écriture inclusive\cite{MEI} de Mots-Clés, s'appuyant sur les travaux de HOUDEBINE Anne-Marie\footnote{Professeure émérite de linguistique et sémiologie Université Paris Descartes – Sorbonne.}.
Ceci permettant de ne plus utiliser le genre masculin comme genre neutre, ce qui exclue 50\% de la population.
Ainsi selon le Manuel d’Écriture inclusive : \og Le point milieu permet d'affirmer sa fonction singulière d'un point de vue sémiotique et par là d'investir "frontalement" l'enjeu discursif et social de l'égalité femmes/hommes. \fg{}
Je vois personnellement, ce rapport comme un laboratoire vivant, ce qui me permettra de statuer au mieux sur l'utilisation d'une écriture égalitaire et novatrice pour la rédaction de l'ensemble de mes écrit.


\newpage


\section*{Résumé}
\addcontentsline{toc}{section}{Résumé}

Dans le but de promouvoir de manière ludique l'engagement associatif, j'ai décidé de rejoutre l'UV AV00\footnote{Unité de valeur hors emploi du temps, ayant pour objectif de réaliser un projet audivisuel.} au cours de mon neuvième semestre d'étude supérieur à l'Université de Technologie de Belfort Montbéliard, pour réaliser un projet audivisuel.

Ce cours documentaire sera esentiellement centré sur les activités estudiantines de l'UTBM, par le biais de la Junior UTBM, de l'AE UTBM et du BDS UTBM.
Je me permets d'inclure le discours d'une étudiante de l'\gls{ENSEEIHT} qui est partie en Colombie pour un voyage humanitaire, qui est aussi Responsable Qualité de la Junior Entreprise N7 Consulting, et qui est membre active de Women7.

\newpage



\tableofcontents
\newpage



\pagestyle{number}

\section{Présentation}\label{sec:presentation}

%! Author = thibaultchausson
%! Date = 23/07/2023

Depuis mes 17 ans, je suis membre du bureau directeur du club Rethel Château Canoë Kayak\footnote{\href{https://rcck-ardennes.fr}{RCCK}, un club sportif et de loisir de kayak dans le Sud-Ardennes.}.
Il m'a paru inconcevable de ne pas m'engager dans les associations estudiantines de l'\gls{UTBM} durant mon cursus d'élève ingénieur.
De ce fait, j'ai rejoint l'Association des Étudiant·e·s de l'\gls{UTBM} au cours du semestre P23 en tant que Responsable des Relations Extérieures, ainsi que la Junior \gls{UTBM} durant ce même semestre, où j'ai eu la charge de suivre des missions en informatique.
À regrets, j'ai décidé d'arrêter mon engagement à l'\gls{AE} pour me concentrer sur mon double diplôme avec le Master ETI. Malgré tout, j'ai décidé de rester au sein de la Junior \gls{UTBM}, où j'occupe le poste de Responsable Qualité durant le semestre A23.
En ce qui me concerne, l'associatif est une expérience humaine très enrichissante, où nous apprenons sans cesse des autres et partageons de forts moments de complicité.

L'association faisant partie intégrante de ma vie et de celle de nombreuses personnes qui m'entourent, c'est en toute logique que j'ai décidé de réaliser mon projet audiovisuel sur ce domaine.
En réalisant ces courtes vidéos montrant la richesse et la diversité de la vie estudiantine, je souhaite la promouvoir et inciter davantage d'étudiant·e·s à la rejoindre dans les prochains semestres.

Ainsi, je vais réaliser, au cours de ce semestre, des vidéos de promotion de cette vie parallèle aux différents impératifs scolaires.
Cette année, j'ai eu l'occasion de suivre de l'intérieur la plus grosse activité de l'Association des Étudiant·e·s, en tant que commis de cuisine pour l'Intégration de l'\gls{UTBM}.
Outre ces trois semaines intenses, la vie associative et l'engagement étudiant ne s'arrêtent jamais : le Bureau Des Sports (BDS) organise des créneaux sportifs quotidiennement, la Junior \gls{UTBM} propose régulièrement des missions rémunérées et professionnalisantes à l'ensemble des étudiant·e·s de l'\gls{UTBM}.
Bien sûr, les domaines des associations ne se limitent pas seulement aux soirées, aux sports, ou à l'entrepreneuriat.

\newpage

\section{Cahier des charges}\label{sec:cahier-des-charges}

%! Author = thibaultchausson
%! Date = 23/07/2023

\subsection{Sujet}\label{subsec:sujet}

Le but de ce projet audiovisual est de réaliser de courtes vidéos appelé Reel sur Instagram, pour montrer la diversité de l'engagement associatif, mais plus particulièrement l'engagement des étudiant·e·s.

Le projet montrera divers moments de vie estudiantine, des interviews d'étudiant·e·s, pour promouvoir les activités dont il·elle·s sont responsable au sein de l'AE UTBM.

Ainsi, l'objectif principal de ce projet est de communiquer sur les différents événements de l'AE, ainsi que ses temps fort, je peux citer de manière non exhaustive, le nouveau principe d'élection pour le Pôle Des Festivités (PDF), les soirées d'ampleurs co-réaliser avec les BDE de Belfort, mais encore promouvoir l'engagement dans l'association\ldots

Comme évoqué procurement, la motivation première de ce sujet est la promotion de la vie associative estudiantine.
Je suis motivé pour promouvoir l'AE en lui donnant une autre image que celle actuelle, en montrant que derrière cette association tentaculaire, il y a des étudiant·e·s qui sont motivés pour améliorer le cadre de vie de leurs paires.
De plus, je trouve que donner la parole aux différents instigateur·trice·s de ces changements est d'autant plus gratifiant pour eux.

\textcolor{red}{Justification du sujet : objectif, motivation (en quoi la forme aide le fond) changement important dopamine}


\subsection{Méthodes}\label{subsec:methodes}


\subsubsection{Le fond}


\subsubsection{La forme}

Pour mener à bien ce projet, j'ai décidé de réaliser de courtes capsules vidéos dynamiques de 1 à 1 minute 30, ce qui permet de garder mon auditoire attentif au message que je souhaite communiquer.
De ce fait, le format de Reels Instagram\footnote{Un Reel est une vidéo courte et divertisante de moins de 90 secondes (\href{https://about.instagram.com/fr-fr/features/reels}{Reels Instagram}).} correspond parfaitement à ces contraintes.




Je peux lister différentes méthodes :
\begin{itemize}
    \item Vidéo d'illustration (avec ou sans bruit du fond)
\end{itemize}

\subsection{Diagramme de Gantt}\label{subsec:diagramme-de-gantt}


\textcolor{red}{Planning initiale  (Gantt) + Commentaire en fin de projet}






\newpage

\section{Réalisation}\label{sec:realisation}

%! Author = thibaultchausson
%! Date = 23/07/2023

\subsection{Formalisation du message}

\textcolor{red}{Formalisation du message, story board}



\subsection{Les informations transmises}

\textcolor{red}{Justification des informations transmise, pourquoi le message, pourquoi ce choix technique}


\subsection{Réalisation des vidéos}





\subsubsection{L'\gls{AE} en Réel}

\subsubsection{Élection du \gls{PDF}}




\subsection{Scripts}\label{subsec:scripts}

\begin{dialogue}
    \speak{Thibault} \lipsum[1]
    \speak{Laura} \lipsum[2]
\end{dialogue}

\newpage


\section{Conclusion}\label{sec:conclusion}




\newpage
\pagenumbering{roman}


\addcontentsline{toc}{section}{Glossaire}
\printglossaries


\newpage

\section*{Annexes}
\addcontentsline{toc}{section}{Annexes}
%! Author = thibaultchausson
%! Date = 23/07/2023



\subsection*{Descriptions des Reels}\label{subsec:descriptions-reels}
\addcontentsline{toc}{subsection}{Descriptions des Reels}



[L'AE en réel ep.1]

\noindent Salut à toi jeune étudiant !

\noindent On démarre cette nouvelle année en te présentant la vie associative de l'UTBM à travers différents réels !
\noindent Au programme, présentation de l'AE, des élections du pôle des festivités et pleins d'autres choses !
\noindent N'hésite pas à venir à notre rencontre, nous poser tes questions ou proposer tes idées pour améliorer la vie étudiante.

\noindent À très vite !

\subsection*{Interface Instagram}\label{subsec:interface-instagram}
\addcontentsline{toc}{subsection}{Interface Instagram}

\begin{figure}[!h]
    \begin{center}
        \includegraphics[scale=0.2]{ressources/interfaceInsta}
        \caption{Interface des Reels Instagram iOS \label{fig:interfaceInsta}}
    \end{center}
\end{figure}



%Formats de fichiers

%Captures


\newpage

\listoffigures
\addcontentsline{toc}{section}{Table des images}


%\newpage

%\listoftables
%\addcontentsline{toc}{section}{Table des tableaux}


\newpage

%\nocite{*} %Pour tout cité, même si non appelé
\bibliography{refs}
\bibliographystyle{plain}
\addcontentsline{toc}{section}{Références}


\end{document}
